\documentclass[11pt]{article}
% ============================================================
%  Namhyoung Kim  |  CV style after the FINX Lab (Prof. Song) CV
%  XeLaTeX, Erewhon (English CV for the homepage)
% ============================================================
\usepackage[a4paper,left=1.5cm,right=1.5cm,top=1.4cm,bottom=1.8cm,footskip=0.9cm]{geometry}
\usepackage{fontspec}
\setmainfont{Erewhon}[Extension=.otf, UprightFont=*-Regular, BoldFont=*-Bold,
  ItalicFont=*-Italic, BoldItalicFont=*-BoldItalic, SlantedFont=*-RegularSlanted]
\usepackage{xcolor}
\definecolor{navy}{rgb}{0.0,0.23,0.44}
\usepackage{enumitem}
\usepackage{lastpage}
\usepackage{fancyhdr}
\usepackage[hidelinks,pdfencoding=auto]{hyperref}

\setlength{\parindent}{0pt}
\setlength{\parskip}{0pt}
\widowpenalty=10000 \clubpenalty=10000
\AtBeginDocument{\raggedright\hyphenpenalty=10000\exhyphenpenalty=50}

\pagestyle{fancy}\fancyhf{}\renewcommand{\headrulewidth}{0pt}
\cfoot{\thepage\ of \pageref{LastPage}}

\makeatletter
\newcommand{\keepspace}[1]{\par\begingroup
  \dimen@=#1\relax\dimen@ii=\pagegoal\advance\dimen@ii-\pagetotal
  \ifdim\dimen@>\dimen@ii\ifdim\dimen@ii>\z@\vfil\fi\break\fi\endgroup}
\makeatother

\newlength{\indA}\setlength{\indA}{1.08cm}   % sub-heading indent
\newlength{\indB}\setlength{\indB}{2.08cm}   % content indent

% ---------- header ----------
\newcommand{\cvheader}{%
  {\color{navy}\huge\bfseries Namhyoung Kim}\par
  \vspace{5pt}\hrule height 0.6pt\par
  \vspace{12pt}
  \begin{tabular}{@{}p{\dimexpr\indB-\tabcolsep\relax}l@{}}
    \textbf{Email}    & \href{mailto:ai.namhyoung@gmail.com}{ai.namhyoung@gmail.com}\\
    \textbf{Phone}    & +82 10-4194-6929\\
    \textbf{Homepage} & \href{https://x7jeon8gi.github.io}{x7jeon8gi.github.io}\\
    \textbf{GitHub}   & \href{https://github.com/x7jeon8gi}{github.com/x7jeon8gi}\\
    \textbf{LinkedIn} & \href{https://www.linkedin.com/in/ai-namhyoung}{linkedin.com/in/ai-namhyoung}\\
    \textbf{Scholar}  & \href{https://scholar.google.com/citations?user=YbxuDSIAAAAJ}{scholar.google.com/citations?user=YbxuDSIAAAAJ}\\
  \end{tabular}\par}

% ---------- sections ----------
\newcommand{\cvsection}[1]{%
  \par\vspace{13pt}\keepspace{72pt}%
  {\color{navy}\Large\bfseries\scshape #1}\par\nobreak\vspace{5pt}}
\newcommand{\cvsub}[2][]{%
  \par\vspace{3pt}\keepspace{56pt}%
  \hspace*{\indA}{\large\scshape #2}\ifx\relax#1\relax\else\ {\footnotesize #1}\fi\par\nobreak\vspace{4pt}}
\newcommand{\corr}{{\footnotesize ($^{\ast}$: Corresponding)}}

% content block indented to the value column
\newenvironment{ind}{\list{}{\leftmargin=\indB\rightmargin=0pt\topsep=0pt\partopsep=0pt\parsep=0pt\itemsep=0pt}\item\relax\def\\{\par\nobreak}}{\endlist}
% organisation line in small caps (inside ind)
\newcommand{\org}[1]{\par\keepspace{48pt}{\large\scshape #1}\par\nobreak}
\newcommand{\gap}{\par\vspace{6pt}}

% numbered list, labels given explicitly: \item[{[4]}]
\newenvironment{nlist}{\begin{list}{}{\leftmargin=2.05cm\labelwidth=0.6cm\labelsep=0.22cm
  \itemsep=5pt\topsep=1pt\partopsep=0pt\parsep=0pt\let\makelabel\nlabel}\interlinepenalty=10000 \def\\{\par\nobreak}}{\end{list}}
\newcommand{\nlabel}[1]{\hfil #1}
% bullets as in the reference
\newenvironment{blist}[1][]{\begin{itemize}[label=\textbullet,leftmargin=2.05cm,labelsep=0.3cm,
  itemsep=5pt,topsep=1pt,partopsep=0pt,parsep=0pt,#1]}{\end{itemize}}
\newenvironment{dlist}{\begin{itemize}[label=\textendash,leftmargin=0.75cm,labelsep=0.3cm,
  itemsep=2pt,topsep=4pt,partopsep=0pt,parsep=0pt]}{\end{itemize}}

\newcommand{\me}{\textbf{Namhyoung Kim}}
\newcommand{\mefirst}{\textbf{Kim, Namhyoung}}
\newcommand{\song}{Jae Wook Song$^{\ast}$}
\newcommand{\lastupdate}{\par\vspace{18pt}\hfill{\small Last Update: 2026-10-13}\par}

\usepackage{kotex}
\setmainhangulfont{NanumMyeongjo}[Path=fonts/, Extension=.ttf,
  UprightFont=*, BoldFont=*Bold, ItalicFont=*, BoldItalicFont=*Bold]
\renewcommand{\corr}{{\footnotesize (\cvmark{*}: 교신저자)}}
\hypersetup{pdftitle={김남형, 이력서},pdfauthor={김남형}}
\renewcommand{\lastupdate}{\par\vspace{18pt}\hfill{\small 최종 갱신: 2026-10-23}\par}

\begin{document}
{\color{navy}\huge\bfseries 김남형\hspace{9pt}{\Large Namhyoung Kim}}\par
\vspace{5pt}\hrule height 0.6pt\par
\vspace{12pt}
\begin{tabular}{@{}p{\dimexpr\indB-\tabcolsep\relax}l@{}}
  \textbf{이메일} & \href{mailto:ai.namhyoung@gmail.com}{ai.namhyoung@gmail.com}\\
  \textbf{홈페이지} & \href{https://x7jeon8gi.github.io/ko/}{x7jeon8gi.github.io/ko/}\\
  \textbf{GitHub} & \href{https://github.com/x7jeon8gi}{github.com/x7jeon8gi}\\
  \textbf{LinkedIn} & \href{https://www.linkedin.com/in/ai-namhyoung}{linkedin.com/in/ai-namhyoung}\\
  \textbf{Scholar} & \href{https://scholar.google.com/citations?user=YbxuDSIAAAAJ}{scholar.google.com/citations?user=YbxuDSIAAAAJ}\\
\end{tabular}\par

% ------------------------------------------------------------
\cvsection{연구 관심 분야}
\begin{ind}
금융 시장을 주된 실증 무대로 삼아 신호에 비해 잡음이 큰 데이터의 표현학습을 연구합니다. 이산 잠재 팩터 모형, 구조 기반 학습, 해석 가능하고 집단에 강건한 모형에 관심이 있습니다.
\end{ind}

% ------------------------------------------------------------
\cvsection{학력}
\begin{blist}
\item 한양대학교 산업공학 박사, 2021.03 – 2026.02
  \begin{dlist}
  \item 학위논문: \textit{Representation Learning in Financial Markets for Portfolio Management}
  \item Financial Innovation \& Decision Analytics (FINX) Lab 연구원
  \item 지도교수: 송재욱
  \item GPA: 4.31 / 4.50
  \end{dlist}
\item 한양대학교 경제금융학 학사 (파이낸스경영 부전공), 2016.03 – 2020.02
  \begin{dlist}
  \item 우등졸업 (Cum Laude)
  \item GPA: 4.00 / 4.50
  \end{dlist}
\end{blist}

% ------------------------------------------------------------
\cvsection{연구 및 경력}
\begin{ind}
\org{Risk X}
Risk X tech Lab 연구소장, 2026.03 – 현재\\
Quantitative AI Researcher, Risk X tech Lab, 2025.07 – 2026.02
\gap
\org{한양대학교}
Financial Innovation \& Decision Analytics (FINX) Lab 연구원, 2021.03 – 2026.02\\
지도교수: 송재욱
\gap
\org{한양대학교}
Business \& A.I. Lab 학부 연구 인턴, 2019.08 – 2019.12\\
지도교수: 김종우
\end{ind}

% ------------------------------------------------------------
\cvsection{연구 실적}
\cvsub[\corr]{국제 학술지 및 학회 논문}
\begin{nlist}
\item[{[4]}] \mefirst, and \song. ``Vector-Quantized Discrete Latent Factors Meet Financial Priors: Dynamic Cross-Sectional Stock Ranking Prediction for Portfolio Construction.'' \textit{Proceedings of the 35th International Joint Conference on Artificial Intelligence} (IJCAI-ECAI 2026). 채택률: 13.6\%.
\item[{[3]}] Na, Yosep, \me, Joon Ho Jun, and \song. ``HAR-CEV: Probabilistic Volatility Forecasting with Adaptive Mean Reversion under Distribution Shift.'' \textit{Proceedings of the 7th ACM International Conference on AI in Finance} (ICAIF '26). 게재 예정.
\item[{[2]}] \mefirst, Yosep Na, and \song. ``Deep Mean-Reversion: A Physics-Informed Contrastive Approach to Pairs Trading.'' \textit{Proceedings of the 6th ACM International Conference on AI in Finance} (ICAIF '25): 405–412. 구두발표, 채택률: 15.5\%.
\item[{[1]}] \mefirst, Seung Eun Ock, and \song. ``FactorVQVAE: Discrete Latent Factor Model via Vector Quantized Variational Autoencoder.'' \textit{Knowledge-Based Systems} 318 (2025): 113460.
\end{nlist}
\cvsub[\corr]{국내 학술지}
\begin{nlist}
\item[{[1]}] 옥승은, 변준영, \textbf{김남형}, 송재욱\cvmark{*}. ``기계학습 기반 시군구 건강수명 산정 모형 및 경상북도 데이터를 이용한 실증.'' \textit{경영과학, 41(1), 51–70}, 2024.
\end{nlist}
\cvsub[\corr]{심사 중 논문 및 워킹페이퍼}
\begin{nlist}
\item[{[3]}] \mefirst, Yosep Na, Jiwon Park, and \song. ``LEXICON: Making Factors Talk with Language-Grounded Anchors in Stock Return Prediction.'' 수정 후 재심사, ACM SIGKDD 2027.
\item[{[2]}] \mefirst, and \song. ``PRAXIS: Policy-Coupled Group Robustness for Interpretable Deep Hedging.''
\item[{[1]}] \mefirst, Seungwon Shin, Jiwon Park, and \song. ``Decision Space Quantization.''
\end{nlist}

% ------------------------------------------------------------
\cvsection{학회 발표}
\cvsub[{\footnotesize (\cvmark{*}: 교신저자, \cvmark{\textdagger}: 공동 제1저자)}]{국제 학술대회}
\begin{nlist}
\item[{[6]}] Park, Jiwon\cvmark{\textdagger}, \me\cvmark{\textdagger}, Seungwon Shin, and \song. ``Text as Structure, Not Signal: Text-Aligned Groups for Distributionally Robust Stock Ranking.'' \textit{Workshop on Interpretability and Alignment of Financial Models, The 7th ACM International Conference on AI in Finance} (IAFM '26), 2026.11, Milan, Italy. 발표 예정.
\item[{[5]}] \mefirst, and \song. ``Vector-Quantized Discrete Latent Factors Meet Financial Priors: Dynamic Cross-Sectional Stock Ranking Prediction for Portfolio Construction.'' \textit{Asian Quantitative Finance Conference} (AQFC 2026), 2026.07, Pohang, Korea.
\item[{[4]}] Na, Yosep, Jun Young Byun, Jungyoon Song, \me, and \song. ``Probabilistic Forecasting of High-Frequency Crypto Volatility via Decomposition-Integrated Recurrent Neural Networks.'' \textit{Workshop on Rethinking Financial Time-Series, The 6th ACM International Conference on AI in Finance} (ICAIF-W), 2025.11, Singapore.
\item[{[3]}] \mefirst, Janghyuk Youn, and \song. ``FPQ-VAE: A Dynamic Factor Model Fusing Financial Priors and Vector-Quantized Factors for Stock Rank Prediction.'' \textit{Workshop on Advances in Financial AI, The 34th ACM International Conference on Information and Knowledge Management} (CIKM-W), 2025.11, Seoul, Korea.
\item[{[2]}] Song, Jae Wook\cvmark{*}, \me, and Dongwon Ryu. ``Deep Hedging of Equity-Linked Securities.'' \textit{INFORMS Annual Meeting}, 2024.10, Seattle, WA, USA.
\item[{[1]}] \mefirst, Seung Eun Ock, and \song. ``Risk Propensity-specific Portfolio Recommendation via Self-supervised Learning.'' \textit{Workshop on Recommender Systems in Finance (Fin-RecSys), The 33rd International Joint Conference on Artificial Intelligence} (IJCAI-W), 2024.07, Jeju, Korea.
\end{nlist}
\cvsub[\corr]{국내 학술대회}
\begin{nlist}
\item[{[7]}] \textbf{김남형}, 송재욱\cvmark{*}. ``Vector-Quantized Discrete Latent Factors Meet Financial Priors: Dynamic Cross-Sectional Stock Ranking Prediction for Portfolio Construction.'' \textit{한국산업경영시스템학회 추계학술대회}, 2026.10, 제주.
\item[{[6]}] \textbf{김남형}, 송재욱\cvmark{*}. ``PRISM-VQ: Fusing Financial Priors and Vector-Quantized Latent Factors for Dynamic Cross-Sectional Stock Prediction.'' \textit{한국파생상품학회 추계학술대회}, 2025.12, 서울.
\item[{[5]}] \textbf{김남형}, 김주환, 송재욱\cvmark{*}. ``Vector Quantization for Efficient Financial Representation Learning.'' \textit{대한산업공학회·한국경영과학회·한국시뮬레이션학회 춘계공동학술대회}, 2025.06, 제주.
\item[{[4]}] \textbf{김남형}, 옥승은, 송재욱\cvmark{*}. ``FactorVQVAE: Discrete Latent Factor Model via Vector Quantized Variational Autoencoder.'' \textit{재무금융 공동학술대회}, 2025.05, 서울.
\item[{[3]}] 김주환, \textbf{김남형}, 옥승은, 송재욱\cvmark{*}. ``Pairs Trading via Contrastive Learning.'' \textit{한국경영과학회·대한산업공학회·한국시뮬레이션학회 춘계공동학술대회}, 2024.05, 여수.
\item[{[2]}] 박민호, \textbf{김남형}, 변준영, 송재욱\cvmark{*}. ``Data Augmentation in Embedding Space for Natural Language Processing based on the Vector Quantized Generative Adversarial Networks.'' \textit{대한산업공학회·한국경영과학회 춘계공동학술대회}, 2023.06, 제주.
\item[{[1]}] \textbf{김남형}, 최영훈, 박민호, 송재욱\cvmark{*}. ``Can a Pre-trained Language Model Be a Superior Counsel in Equity Research Report-Driven Investment Strategy?.'' \textit{한국경영과학회·대한산업공학회 춘계공동학술대회}, 2022.06, 제주.
\end{nlist}

% ------------------------------------------------------------
\cvsection{연구과제 및 산학 프로젝트}
\cvsub{정부 연구과제}
\begin{nlist}
\item[{[6]}] 신경망 기반 최적제어를 활용한 ELS 헤징 의사결정 시스템 개발\\ 서울경제진흥원(SBA), 핀테크 기술사업화 지원사업, 2026.07 – 현재. (과제 총괄)
\item[{[5]}] 시계열 생성 기반 투자자 맞춤 구조화 파생상품 개발\\ 팁스(TIPS) R\&D, 2026.04 – 현재. (기술 총괄)
\item[{[4]}] 벡터 양자화 변분 오토인코더를 활용한 금융 데이터의 강건한 생성적 표현 학습 및 응용\\ 한국연구재단, 2025.09 – 2026.08.
\item[{[3]}] 벡터 양자화 변분 오토인코더를 활용한 포트폴리오 분야 인공지능 응용연구\\ 한국연구재단, 2023.06 – 2025.06. (연구원)
\item[{[2]}] 설명 가능한 인공지능을 활용한 대출자격 분류모형 연구\\ 한국연구재단, 2021.06 – 2023.02. (연구원)
\item[{[1]}] 공유와 협력에 기반한 스마트 시스템 교육연구팀\\ 한국연구재단, 4단계 BK21, 2021.03 – 2025.07.
\end{nlist}
\cvsub{산학 프로젝트}
\begin{nlist}
\item[{[7]}] AI 기술을 활용한 금융 활동 분석\\ 금융결제원(KFTC), 2025.03 – 2025.07. (연구원, 모형 설계 및 구현)
\item[{[6]}] 퇴직연금 관리를 위한 기계학습 기반 모델 포트폴리오\\ FnDataLab, 2024.03 – 2024.06. (연구원)
\item[{[5]}] 기계학습을 통한 이색 옵션 및 구조화상품 프라이싱\\ Risk X tech Lab, 2023.10 – 2024.01. (연구원)
\item[{[4]}] 기계학습 기반 글로벌 투자자산 위험등급 산정\\ FnDataLab, 2023.05 – 2023.08. (연구원, 모형 설계 및 구현)
\item[{[3]}] 그래프 신경망과 자연어 처리 기반 세일즈 애널리틱스\\ IOI Soft, 2022.11. (연구원)
\item[{[2]}] 중소기업 맞춤형 비즈니스 애널리틱스 개발\\ 하이퍼라운지, 2021.10 – 2021.11. (연구원)
\item[{[1]}] 미세먼지 저감 산림사업 효과성 분석\\ 한국임업진흥원(KOFPI), 2021.05 – 2021.10.
\end{nlist}

% ------------------------------------------------------------
\cvsection{특허}
\begin{blist}
\item 국내 특허 3건 출원, 2026: 리 대수 기반 시장 시나리오 생성, 구조화상품 딥헤징 시스템 등.
\end{blist}

% ------------------------------------------------------------
\clearpage
\cvsection{수상 및 장학금}
\begin{blist}[itemsep=3pt]
\item 최우수상 (박사 트랙), 한국산업경영시스템학회 제1회 대학원 우수 학술연구 발표대회, 2026.10
\item 우수논문상, 한국경영과학회, 2024.11
\item BK21 Four 연구장학금, 한국연구재단, 2021.03 – 2025.07
\item AI 대학원 지원사업 장학금, 정보통신기획평가원(IITP), 2020.09 – 2020.10
\item 우등졸업 (Cum Laude), 한양대학교, 2020.02
\item 금융 퀀트 전문 양성 사업 장학금, 한양대학교, 2018.12
\end{blist}

% ------------------------------------------------------------
\cvsection{교육 경력}
\cvsub{한양대학교}
\begin{blist}
\item 공학도를 위한 창의적 프로그래밍(11391) 조교, 한양대학교, 2025년 1학기.
\end{blist}

% ------------------------------------------------------------
\cvsection{기술}
\begin{blist}
\item {프로그래밍 및 도구}: Python, PyTorch, \LaTeX, Git, Docker, Linux, Weights \& Biases
\end{blist}

% ------------------------------------------------------------
\cvsection{추천인}
\begin{ind}
송재욱\\
한양대학교 산업공학과 부교수\\
\href{mailto:jwsong@hanyang.ac.kr}{jwsong@hanyang.ac.kr}
\end{ind}

\lastupdate
\end{document}
